\documentclass[
    language=english,
    doctype=sermon,
    institution=none,
]{../../omnilatex}

\RequirePackage{config/document-settings}

\begin{document}

\maketitle

\section{The Parable of the Lost Coin}

\textbf{Scripture Reading:} Luke 15:8--10 (ESV)

\begin{quote}
``Or what woman, having ten silver coins, if she loses one coin, does not light a lamp and sweep the house and seek diligently until she finds it? And when she has found it, she calls together her friends and neighbors, saying, `Rejoice with me, for I have found the coin that I had lost.' Just so, I tell you, there is joy before the angels of God over one sinner who repents.''
\end{quote}

\section{Introduction}

The parable of the lost coin is the second of three parables in Luke chapter 15---sometimes called ``The Shepherd's Heart,'' ``The Woman's Persistence,'' and ``The Father's Joy.'' Together, these three stories reveal a God who is actively, personally, and passionately searching for what has been lost.

Today we focus on the middle parable: a woman who loses one of ten silver coins. This was not pocket change. In first-century Palestine, these coins were likely part of a woman's dowry---her financial security, her identity, her life savings. To lose one coin was a genuine crisis.

Let us consider three truths from this parable.

\section{Point 1: The Value of the Lost}

\textbf{What was lost?} A silver coin---a drachma, worth approximately one day's wages for a laborer. Ten coins represented a woman's entire savings, her security, her connection to her community.

When we read this parable, we are invited to see ourselves in the coin. Each of us matters to God. Not in an abstract, theological sense, but in a personal, urgent way. The woman did not say, ``I still have nine; one missing is acceptable.'' She said, ``I have lost one,'' and she lit a lamp.

This is the heart of the gospel: \textbf{you are not a statistic to God.} You are the one. The one who wandered. The one who is hiding. The one who feels forgotten. God does not move on. He lights a lamp.

\bigskip
\noindent\textit{Application:} If you are here today and you feel lost---perhaps you have drifted from faith, perhaps you are carrying a burden that makes you feel unworthy---know that you are the one God is searching for. Your value is not determined by your productivity, your moral record, or your usefulness. You are valued because you are His.

\section{Point 2: The Search}

\textbf{How did she search?} She lit a lamp. She swept the house. She searched diligently.

This was a thorough, systematic, physical search. She did not sit down and wish the coin would return. She got up, lit a lamp (oil was expensive), and swept every corner. The Greek word \emph{diakopōs} (translated ``diligently'') carries the sense of ``with intensity'' or ``without ceasing.''

This is a picture of God's pursuit. He does not passively wait. He actively seeks. Jesus told these parables because the Pharisees were grumbling that He ``welcomes sinners and eats with them'' (Luke 15:2). Jesus' response is devastating in its simplicity: \textbf{God is not a critic of the lost; He is a searcher of the lost.}

\bigskip
\noindent\textit{Application:} Are you searching for someone---a prodigal child, a friend who has walked away, a colleague who seems unreachable? Do not give up. The woman searched every corner. Some of the most important ministry we do is the quiet, patient, persistent search for one person who matters to God.

\section{Point 3: The Celebration}

\textbf{What happened when she found it?} She called her friends and neighbors. She threw a party. ``Rejoice with me, for I have found the coin that I had lost.''

Notice the invitation: \textit{Rejoice with me.} This was not a private celebration. The joy was communal. In Jewish culture, when something precious was recovered, the community gathered to share in the relief and delight.

Jesus draws the parallel directly: ``Just so, I tell you, there is joy before the angels of God over one sinner who repents'' (Luke 15:10). When one person turns back to God, all of heaven celebrates. The angels---beings who have never sinned---rejoice over the restoration of one broken human being.

This is the scandal of grace: \textbf{the party is for the one, not for the many.} We tend to celebrate big events---large gatherings, baptisms of crowds, churches that grow. But heaven's celebration is tuned to the frequency of the individual. One coin. One sheep. One son. One sinner who comes home.

\bigskip
\noindent\textit{Application:} Do you know the joy of having been found? Have you experienced the moment when the weight of your sin was lifted, when the God who searched for you finally found you and threw a party? If you have, then let that memory fuel your own search for the lost around you. If you have not, today may be the day the lamp is lit for you.

\section{Closing Prayer}

\begin{quote}
Heavenly Father,

We thank You that You are not a distant, indifferent God, but a searching, pursuing, rejoicing God. Thank You that when we were lost, You did not give up on us. Thank You for the light of Your Word, the sweep of Your Spirit, and the diligence of Your love.

For those among us who feel lost today---who carry the weight of sin, shame, or distance---we ask that You would find them. Let them know they are the one. Let them feel the warmth of the lamp you have lit. Let them hear the joy of heaven.

And for all of us, Lord, give us the heart of the woman who searched. Give us eyes to see the lost around us, courage to light the lamp, and joy to celebrate when they are found.

In the name of Jesus Christ, who came to seek and save the lost.

Amen.
\end{quote}

\section{Benediction}

Go now in the peace of Christ. May the grace of the Lord Jesus, the love of God, and the fellowship of the Holy Spirit be with you all. And may you carry with you the joy of the one who was found.

\textbf{Amen.}

\end{document}
